\section{Introducción}

Textiles y Confecciones Atlas E.I.R.L. es una empresa con más de 35 años de prestigiosa trayectoria en la ciudad de Arequipa, dedicada a la confección y comercialización de indumentaria deportiva, uniformes institucionales y artículos vinculados para clubes, colegios y empresas mineras. A pesar de su sólida presencia física, la compañía operaba bajo un modelo de madurez digital inicial, con una dependencia crítica de procesos manuales que limitaban severamente su escalabilidad y alcance regional.

Antes de la intervención tecnológica, la operación comercial se enfrentaba a múltiples dolores críticos. En primer lugar, existía un evidente "Punto Único de Falla", ya que toda la atención dependía de una sola computadora en el mostrador y un único celular corporativo, lo que paralizaba la facturación ante cualquier avería técnica. Asimismo, la facturación se realizaba de manera aislada en el portal web de la SUNAT, obligando a los vendedores a realizar una doble digitación que generaba un desfase constante entre lo vendido y el stock físico (quiebres de stock fantasma). A nivel estratégico, la falta de datos centralizados obligaba a la gerencia a dedicar hasta el 70\% de su tiempo en tareas netamente operativas, como la verificación física en los estantes. Adicionalmente, la empresa sufría de un "Silencio Digital" con nula presencia oficial en redes sociales, y presentaba una alta vulnerabilidad informática al respaldar su información precariamente en memorias USB, sin políticas de ciberseguridad.

Ante este desafiante escenario, entre fines de octubre de 2025 y julio de 2026, y bajo el marco de cofinanciamiento con ProInnóvate, se ejecutó un proceso integral de transformación digital estructurado en un horizonte de 6 meses. La solución central consistió en la implementación del sistema ERP Odoo en su versión SaaS (Cloud) ---digitalizando el inventario y modernizando el Punto de Venta (POS)--- acompañado de una profesionalización en marketing omnicanal.

El presente informe tiene como finalidad evaluar los resultados de esta implementación, analizando el impacto de la digitalización en la operatividad de la empresa, la optimización de los tiempos de atención y el establecimiento de canales digitales que modernizaron la propuesta comercial de la organización.
